\documentclass[12pt,a4paper]{article}

\usepackage[a4paper,margin=1in]{geometry}
\usepackage{graphicx}
\usepackage{booktabs}
\usepackage{amsmath,amssymb}
\usepackage{float}
\usepackage{hyperref}
\usepackage{xcolor}
\usepackage{listings}
\usepackage{enumitem}
\usepackage{fancyhdr}
\usepackage{longtable}

\hypersetup{colorlinks=true,linkcolor=blue,urlcolor=blue}

\pagestyle{fancy}
\fancyhf{}
\lhead{ICS1512}
\rhead{Machine Learning Algorithms Laboratory}
\cfoot{\thepage}

\lstset{
language=Python,
basicstyle=\ttfamily\small,
numbers=left,
frame=single,
breaklines=true,
showstringspaces=false
}

\begin{document}

\begin{center}
{\large \textbf{Sri Sivasubramaniya Nadar College of Engineering, Chennai}}\\
{\small (An Autonomous Institution affiliated to Anna University)}
\end{center}
\vspace{0.3cm}

\begin{tabular}{|p{4cm}|p{11cm}|}
\hline
Experiment No. & 9 \\ \hline
Experiment Title & Perceptron vs Multilayer Perceptron (A/B Experiment) with Hyperparameter Tuning \\ \hline
Student Name & L.Radesh \\ \hline
Register Number & 3122247001047 \\ \hline
Faculty & Dr. Poreddy Ajay Kumar Reddy \\ \hline
Submission Date & 21/09/2026 \\ \hline
Github Repo & \url{https://github.com/RadeshL/Machine-Learning-Laboratory/blob/main/MLassign9.ipynb} \\ \hline
\end{tabular}

\section{Aim and Objective}
To implement and compare the performance of two classifiers on the English Handwritten Characters
dataset (3{,}410 images, 62 classes: digits 0--9, uppercase A--Z, lowercase a--z):
\begin{itemize}
\item \textbf{Model A}: a single-layer Perceptron Learning Algorithm (PLA), implemented from scratch.
\item \textbf{Model B}: a Multilayer Perceptron (MLP) with hidden layers and non-linear activations,
with hyperparameters (activation function, optimizer, learning rate, network depth, batch size,
regularization) selected through systematic tuning.
\end{itemize}
The experiment evaluates both models with standard classification metrics and studies how sensitive
each model is to its hyperparameters and to the linear-separability limitations of a single-layer
network.

\section{Dataset Description}
\begin{longtable}{|p{4cm}|p{5cm}|p{1.6cm}|p{1.8cm}|p{2.2cm}|}
\hline
\textbf{Dataset} & \textbf{Source} & \textbf{Images} & \textbf{Classes} & \textbf{Images/class} \\ \hline
English Handwritten Characters & Kaggle (dhruvildave/english-handwritten-characters-dataset) & 3410 & 62 (0--9, A--Z, a--z) & 55 \\ \hline
\end{longtable}

Each class (0--9, A--Z, a--z) has exactly 55 handwritten sample images, giving a perfectly balanced
3{,}410-image dataset. The raw images are $1200\times900$ RGB scans of a single handwritten character
each, with dark ink strokes on a white background.

\section{Preprocessing Steps}
\begin{itemize}
\item Each image was converted to grayscale and resized to $32\times32$ pixels using Lanczos
resampling, giving a manageable 1024-dimensional feature vector per image.
\item Pixel intensities were normalized to $[0,1]$ and then \textbf{inverted} (background $\to 0$,
ink $\to 1$) so that the informative strokes correspond to non-zero activations, which is the
standard convention for this kind of character data and helps both the perceptron and the MLP learn
faster.
\item Labels (0--9, A--Z, a--z) were encoded as integer class indices $0$--$61$ in a fixed,
reproducible sorted order.
\item An 80/20 stratified train/test split (2728 train / 682 test images) was used, keeping the class
balance identical in both splits. For hyperparameter tuning, a further stratified 85/15 split of the
training data (2318 tuning-train / 410 tuning-validation) was used, so that the held-out 682-image
test set is touched only once, at the very end, by the final chosen models.
\end{itemize}

\section{Model A: PLA Implementation and Results}

\subsection{Implementation}
Because this is a 62-class problem, the PLA was implemented as the standard multiclass generalisation
of the perceptron (a weight matrix with one row per class, sometimes called Kesler's construction),
built entirely from scratch in NumPy:
\begin{itemize}
\item \textbf{Step activation / prediction:} $\hat{y} = \arg\max_c \; (W_c \cdot x)$, i.e. the class
whose weight vector gives the highest dot product with the input.
\item \textbf{Weight update rule} (applied only when a sample is misclassified), which is the
multiclass form of $w_{t+1} = w_t + \eta (y - \hat y) x$:
\[
W_{y_{\text{true}}} \mathrel{+}= \eta \, x, \qquad
W_{\hat y} \mathrel{-}= \eta \, x
\]
i.e.\ the true class's weight vector is pushed towards $x$ and the (wrongly) predicted class's weight
vector is pushed away from $x$.
\end{itemize}
Training ran for up to 40 epochs (full passes over the shuffled 2728-image training set) with a
learning rate $\eta = 0.1$, stopping early only if an entire epoch produced zero updates
(i.e.\ the training data became linearly separable under the current weights).

\subsection{Results}
\begin{figure}[H]
\centering
\includegraphics[width=\linewidth]{outputs/pla_convergence.png}
\caption{Left: PLA training/test error vs epoch ($\eta=0.1$). Right: effect of learning rate on the
training-error convergence curve ($\eta \in \{0.01, 0.1, 0.5\}$).}
\end{figure}

\begin{table}[H]
\centering
\caption{PLA -- Final Test-Set Metrics (682 held-out images)}
\begin{tabular}{lcccc}
\toprule
Metric & Accuracy & Precision (macro) & Recall (macro) & F1 (macro) \\
\midrule
PLA & 0.2434 & 0.2478 & 0.2434 & 0.2243 \\
\bottomrule
\end{tabular}
\end{table}

The PLA's training accuracy climbed steadily to about 75\% by epoch 40 (it never fully converges --
the data is not linearly separable), while its test accuracy stayed essentially flat around 20--25\%
throughout training. Varying the learning rate ($\eta \in \{0.01, 0.1, 0.5\}$) changed how quickly the
training error dropped in the first few epochs, but all three settings plateaued at a similar training
error, confirming the limitation is structural (linear boundaries), not a learning-rate problem.

\section{Model B: MLP Implementation and Hyperparameter Tuning}

\subsection{Implementation}
The MLP was built using \texttt{sklearn.neural\_network.MLPClassifier}, a fully-connected feed-forward
network trained with backpropagation and a multi-class cross-entropy loss. Hyperparameters tuned:
activation function, optimizer, learning rate, hidden-layer depth/width, batch size, and L2
regularization strength ($\alpha$).

\subsection{Hyperparameter Sweep}
Eleven configurations (60 epochs each, on the 2318/410 tuning train/validation split) were compared,
sweeping one axis at a time from a ReLU + Adam + one 128-unit hidden layer baseline:

\begin{table}[H]
\centering
\small
\caption{MLP Hyperparameter Sweep Results}
\begin{tabular}{lllrrrr}
\toprule
Config & Activation & Solver & LR & Batch & Train Acc & Val Acc \\
\midrule
A1 (128) & relu & adam & 0.001 & 32 & 0.9957 & 0.4317 \\
A2 (128) & tanh & adam & 0.001 & 32 & 0.9871 & 0.3634 \\
A3 (128) & logistic & adam & 0.001 & 32 & 0.8395 & 0.4390 \\
B1 (128) & relu & sgd & 0.001 & 32 & 0.4888 & 0.3512 \\
B2 (128) & relu & adam & 0.01 & 32 & 1.0000 & 0.4488 \\
B3 (128) & relu & adam & 0.0001 & 32 & 0.5833 & 0.3683 \\
C1 (128) & relu & adam & 0.001 & 16 & 0.9987 & 0.4341 \\
C2 (128) & relu & adam & 0.001 & 64 & 0.9724 & 0.4268 \\
D1 (256,128) & relu & adam & 0.001 & 32 & 0.9978 & 0.4878 \\
\textbf{D2 (256,128,64)} & relu & adam & 0.001 & 32 & 1.0000 & \textbf{0.5146} \\
D3 (512) & relu & adam & 0.001 & 32 & 1.0000 & 0.4707 \\
\bottomrule
\end{tabular}
\end{table}

\begin{figure}[H]
\centering
\includegraphics[width=\linewidth]{outputs/mlp_tuning_barplot.png}
\caption{Validation accuracy across the 11-configuration hyperparameter sweep.}
\end{figure}

\textbf{Reading the sweep:}
\begin{itemize}
\item \textbf{Activation} (A1 vs A2 vs A3): ReLU (A1) beat tanh and logistic -- it avoids the
vanishing-gradient squashing that hurts the other two on this fairly deep, high-dimensional input.
\item \textbf{Optimizer / learning rate} (B1--B3 vs A1): plain SGD (B1) is clearly worse than Adam at
the same learning rate -- Adam's per-parameter adaptive steps converge much faster within the fixed
60-epoch budget. A learning rate that is too high (B2, 0.01) or too low (B3, 0.0001) both underperform
the middle setting (0.001).
\item \textbf{Batch size} (C1 vs C2 vs A1): the mid-sized batch (32) gave the best accuracy/time
trade-off.
\item \textbf{Depth/width} (D1--D3): going from 1 to 2 to 3 hidden layers steadily improved validation
accuracy (D2, 3 layers, is the best configuration overall), while simply widening a single layer (D3)
helped less than adding depth.
\end{itemize}

\subsection{Overfitting Check and Regularization}
The best architecture from the sweep (D2, 3 hidden layers) reaches essentially 100\% training accuracy
but only about 51\% validation accuracy -- a clear sign of overfitting, since the network has more
than enough capacity to memorise the $\sim$2{,}300 training images. An L2-regularization ($\alpha$)
sweep on this architecture was used to address it:

\begin{table}[H]
\centering
\caption{Effect of L2 Regularization (3-hidden-layer MLP)}
\begin{tabular}{lccc}
\toprule
$\alpha$ & Train Acc & Val Acc & Final Loss \\
\midrule
0.001 & 1.0000 & 0.5512 & 0.0284 \\
\textbf{0.01} & 0.9974 & \textbf{0.5585} & 0.1622 \\
0.1 & 0.9948 & 0.5195 & 0.9056 \\
\bottomrule
\end{tabular}
\end{table}

\begin{figure}[H]
\centering
\includegraphics[width=0.7\linewidth]{outputs/mlp_regularization.png}
\caption{Train vs validation accuracy as L2 regularization strength ($\alpha$) increases.}
\end{figure}

$\alpha = 0.01$ gave the best validation accuracy while only slightly trimming training accuracy,
confirming genuine overfitting was present and that weight-decay-style regularization helps close the
gap.

\subsection{Chosen Final Hyperparameters}
\begin{table}[H]
\centering
\caption{Final MLP Configuration}
\begin{tabular}{ll}
\toprule
Hyperparameter & Chosen value \\
\midrule
Hidden layers & (256, 128, 64) -- 3 layers \\
Activation & ReLU \\
Optimizer & Adam \\
Learning rate & 0.001 \\
Batch size & 32 \\
L2 regularization ($\alpha$) & 0.01 \\
Early stopping & on, patience 20 epochs, 15\% internal validation split \\
\bottomrule
\end{tabular}
\end{table}

\subsection{Final Training and Results}
\begin{figure}[H]
\centering
\includegraphics[width=\linewidth]{outputs/mlp_training_curve.png}
\caption{Left: MLP training loss (cross-entropy) vs epoch. Right: internal validation accuracy vs
epoch, with early stopping triggered after 53 epochs (best validation accuracy 0.5244).}
\end{figure}

Training with early stopping enabled halted after 53 epochs once validation accuracy stopped
improving for 20 consecutive epochs, avoiding further overfitting beyond that point.

\begin{table}[H]
\centering
\caption{MLP -- Final Test-Set Metrics (682 held-out images)}
\begin{tabular}{lcccc}
\toprule
Metric & Accuracy & Precision (macro) & Recall (macro) & F1 (macro) \\
\midrule
MLP (tuned) & 0.5132 & 0.5269 & 0.5132 & 0.5075 \\
\bottomrule
\end{tabular}
\end{table}

\section{A/B Comparison: PLA vs Tuned MLP}

\begin{table}[H]
\centering
\caption{Final Comparison on the Held-Out Test Set}
\begin{tabular}{lcccccc}
\toprule
Model & Accuracy & Prec.\ (macro) & Rec.\ (macro) & F1 (macro) & ROC-AUC (micro) & ROC-AUC (macro) \\
\midrule
PLA (single-layer) & 0.2434 & 0.2478 & 0.2434 & 0.2243 & 0.835 & 0.837 \\
MLP (tuned) & \textbf{0.5132} & \textbf{0.5269} & \textbf{0.5132} & \textbf{0.5075} & \textbf{0.961} & \textbf{0.961} \\
\bottomrule
\end{tabular}
\end{table}

\begin{figure}[H]
\centering
\includegraphics[width=0.75\linewidth]{outputs/final_metrics_barplot.png}
\caption{PLA vs tuned MLP on accuracy, macro precision, macro recall and macro F1.}
\end{figure}

The tuned MLP more than doubles the PLA's test accuracy (51.3\% vs 24.3\%) and improves every other
metric by a similarly large margin. This gap is the direct, expected consequence of the two models'
representational power: the PLA can only separate classes with straight-line boundaries in the
1024-dimensional pixel space, while the MLP's hidden layers and ReLU non-linearities let it build much
more complex decision regions.

\textbf{Strengths and weaknesses:}
\begin{itemize}
\item \textbf{PLA}: simple, fast to train (a single epoch touches every sample only once with a
lightweight update), and easy to interpret, but fundamentally limited to linear decision boundaries --
inadequate for 62 visually overlapping handwritten classes.
\item \textbf{MLP}: much higher accuracy thanks to non-linear, multi-layer representations, but slower
to train, needs careful hyperparameter tuning to avoid over/under-fitting, and its decision boundaries
are far less interpretable.
\end{itemize}

\textbf{Impact of hyperparameter tuning:} the untuned baseline MLP configuration (A1: 1 hidden layer,
otherwise similar settings) reached 43.2\% validation accuracy; systematic tuning of depth and
regularization pushed the final held-out test accuracy to 51.3\%, an 8-point improvement purely from
architecture and regularization choices, without adding any extra data.

\section{Confusion Matrices and ROC Curves}

\begin{figure}[H]
\centering
\includegraphics[width=\linewidth]{outputs/confusion_matrices.png}
\caption{Confusion matrices on the 682-image test set (62 classes each). The MLP's matrix is visibly
more concentrated on the diagonal than the PLA's.}
\end{figure}

\begin{figure}[H]
\centering
\includegraphics[width=0.8\linewidth]{outputs/hardest_classes.png}
\caption{The 10 hardest classes for the tuned MLP (lowest per-class recall), compared against the PLA
on the same classes.}
\end{figure}

The hardest classes for the MLP are dominated by characters with strong visual look-alikes across the
digit/uppercase/lowercase groups: \texttt{o}/\texttt{O}/\texttt{0}, \texttt{s}/\texttt{S}/\texttt{5},
\texttt{z}/\texttt{Z}, \texttt{x}/\texttt{X}, \texttt{I}/\texttt{l}/\texttt{1}, and similar pairs --
exactly the confusions a human reader would also find genuinely ambiguous in isolated handwritten
characters with no surrounding context.

\begin{figure}[H]
\centering
\includegraphics[width=0.7\linewidth]{outputs/roc_curves.png}
\caption{One-vs-rest ROC curves (micro- and macro-averaged over all 62 classes) for PLA and the tuned
MLP.}
\end{figure}

The MLP's ROC curves sit far closer to the top-left corner than the PLA's, with micro/macro AUC of
about 0.96 versus 0.84 for the PLA -- consistent with the accuracy and F1 gaps above, and showing that
the MLP's advantage holds not just at a single operating threshold but across the full range of
decision thresholds.

\section{Observations and Analysis}
\begin{itemize}
\item \textbf{Why does PLA underperform compared to MLP?} The PLA can only draw straight-line
(linear) decision boundaries in pixel space. With 62 visually overlapping handwritten classes
(e.g.\ \texttt{o}/\texttt{O}/\texttt{0}, \texttt{l}/\texttt{I}/\texttt{1}), the data is not linearly
separable, so the PLA plateaus at a low test accuracy while the MLP's hidden layers and non-linear
ReLU activations let it carve out far more complex decision regions and roughly double the test
accuracy.
\item \textbf{Which hyperparameters had the most impact on MLP performance?} Network depth
(1$\to$3 hidden layers) and the choice of optimizer (Adam vs plain SGD) had the largest effect on
validation accuracy; activation function and learning rate were the next most influential, while
batch size gave smaller, second-order changes.
\item \textbf{Did optimizer choice (SGD vs Adam) affect convergence?} Yes, substantially -- within
the same 60-epoch budget, Adam converged to a much lower loss and higher accuracy (43.2\% val.\ acc.)
than plain SGD (35.1\% val.\ acc.), which was still climbing slowly. Adam's per-parameter adaptive
learning rates help a lot on this fairly high-dimensional (1024-feature) input.
\item \textbf{Did adding more hidden layers always improve results?} Mostly yes, up to a point:
1$\to$2$\to$3 layers steadily improved validation accuracy (43.2\% $\to$ 48.8\% $\to$ 51.5\%), but
capacity and the training/validation accuracy gap (overfitting risk) also grow with depth, so depth
has to be paired with regularization/early stopping rather than increased without limit -- a very
wide single layer (512 units) under-performed the 3-layer network despite having more parameters.
\item \textbf{Did MLP show overfitting? How could it be mitigated?} Yes, clearly: the unregularized
3-layer MLP reached essentially 100\% training accuracy but only about 51--52\% validation accuracy.
Adding L2 regularization ($\alpha=0.01$) and using early stopping on a held-out validation split both
helped narrow this gap and were used in the final model. Other options not explored here (data
augmentation, dropout, a smaller network) could reduce it further.
\end{itemize}

\section{Discussion}
The central theme of this experiment is that model capacity has to match problem complexity. A
62-class handwritten-character recognition problem, where several classes are visually near-identical,
is not linearly separable, so a single-layer perceptron -- however carefully its learning rate is
tuned -- cannot exceed roughly a quarter of the accuracy that a modestly-sized, properly-regularized
multilayer perceptron achieves. At the same time, the MLP's much larger capacity comes with a real
overfitting risk that has to be actively managed through regularization, early stopping, and a
validation-driven hyperparameter search, rather than assumed away.

\section{Conclusion}
Across this A/B experiment on the Kaggle English Handwritten Characters dataset, the tuned MLP
(3 hidden layers, ReLU, Adam, learning rate 0.001, batch size 32, L2 $\alpha=0.01$, early stopping)
substantially outperformed the from-scratch single-layer Perceptron on every metric evaluated --
accuracy (51.3\% vs 24.3\%), macro/weighted precision, recall and F1, and ROC-AUC (0.96 vs 0.84) --
confirming that a non-linear, multi-layer representation is necessary to separate 62 visually similar
handwritten character classes, while the linear PLA is fundamentally limited by its single decision
boundary per class. Systematic hyperparameter tuning (activation, optimizer, learning rate, depth,
batch size, regularization) was essential to reach the MLP's final accuracy, improving it by roughly
8 percentage points over an untuned baseline configuration.

\section{Learning Outcomes}
\begin{itemize}
\item Implemented the multiclass perceptron learning algorithm and its weight-update rule from
scratch, and observed first-hand its inability to solve a non-linearly-separable, 62-class problem.
\item Practiced systematic hyperparameter tuning of a neural network across activation function,
optimizer, learning rate, batch size, depth/width, and regularization strength, and learned to read
train-vs-validation accuracy gaps as a diagnostic for overfitting.
\item Directly observed the effect of the optimizer (Adam vs SGD) and of L2 regularization / early
stopping on training dynamics and generalization.
\item Practiced computing and interpreting macro and weighted precision/recall/F1, confusion
matrices, and one-vs-rest micro/macro ROC-AUC for a genuinely multi-class (62-class) classification
problem.
\item Understood, through per-class recall analysis, that a model's hardest classes often line up
with humanly-ambiguous look-alike characters, not with arbitrary weaknesses of the classifier.
\end{itemize}

\section{References}
\begin{enumerate}
\item Kaggle: English Handwritten Characters Dataset --
\url{https://www.kaggle.com/datasets/dhruvildave/english-handwritten-characters-dataset}
\item Scikit-learn documentation, \texttt{MLPClassifier}: \url{https://scikit-learn.org/stable/modules/generated/sklearn.neural_network.MLPClassifier.html}
\item Rosenblatt, F., ``The Perceptron: A Probabilistic Model for Information Storage and
Organization in the Brain'', Psychological Review, 1958.
\item Rumelhart, D.E., Hinton, G.E., Williams, R.J., ``Learning representations by back-propagating
errors'', Nature, 1986.
\end{enumerate}

\end{document}